% ============================================================
% AVANCE DE IMPLEMENTACIÓN — SEMANA DEL 21 DE SEPTIEMBRE
% Este bloque está preparado para colocarse después del diseño.
% ============================================================

\section{Implementación del modelo}

La implementación del modelo se desarrolla de manera incremental a partir de
la Especificación de Diseño del Modelo de Evaluación Automatizada de Calidad
Web -- Versión 1.0. En esta etapa se operacionalizan progresivamente la
estructura común de información, los catálogos, las reglas de procesamiento y
la generación del diagnóstico, manteniendo la correspondencia con los
requisitos y elementos de diseño definidos previamente.

El primer incremento se concentra en la base técnica del prototipo. Por esta
razón, todavía no incorpora la ejecución de WAVE, axe DevTools o Lighthouse
sobre el Portal Académico ni presenta resultados del caso de estudio. Su
propósito es establecer una estructura ejecutable y verificable que permita
representar de forma homogénea los datos que posteriormente serán obtenidos de
las tres herramientas.

\subsection{Alcance del primer incremento de implementación}

El alcance de este incremento comprende la preparación del entorno, la
organización interna del proyecto, la implementación de la estructura común
correspondiente a los campos CR-01 a CR-18, la construcción de los catálogos
iniciales y la validación de datos controlados. Las actividades desarrolladas
se presentan en el Cuadro~\ref{tab:actividades-primer-incremento}.

\begin{xltabular}{\textwidth}{p{1.3cm}p{4.2cm}X}
\caption{Actividades del primer incremento de implementación}
\label{tab:actividades-primer-incremento}\\
\toprule
\textbf{ID} & \textbf{Actividad} & \textbf{Resultado obtenido}\\
\midrule
\endfirsthead
\multicolumn{3}{c}{\tablename\ \thetable\ (continuación)}\\
\toprule
\textbf{ID} & \textbf{Actividad} & \textbf{Resultado obtenido}\\
\midrule
\endhead
IM-01 & Preparación del entorno & Proyecto Python ejecutable y dependencias declaradas.\\
IM-02 & Organización del código & Separación entre dominio, catálogos, validación, datos de prueba y pruebas automáticas.\\
IM-03 & Implementación de la estructura común & Modelo tipado para representar los campos CR-01 a CR-18.\\
IM-04 & Implementación de catálogos & Archivos configurables para dimensiones, indicadores, estados, prioridades y reglas.\\
IM-05 & Preparación de datos controlados & Ejemplos representativos de WAVE, axe DevTools y Lighthouse.\\
IM-06 & Verificación inicial & Pruebas de campos obligatorios, enumeraciones y coherencia entre indicador y dimensión.\\
\bottomrule
\end{xltabular}
\begin{center}{\footnotesize\textit{Fuente: Elaboración propia.}}\end{center}

\subsection{Entorno de implementación}

El núcleo del prototipo fue desarrollado en Python debido a que permite
representar estructuras de datos heterogéneas y facilita la posterior
implementación de reglas de transformación, consolidación y análisis. Los
resultados se intercambian mediante archivos JSON, lo cual permite conservar
la estructura original de cada herramienta junto con la representación común
utilizada por el modelo.

La validación estructural se implementó mediante modelos tipados. Esta decisión
permite detectar campos ausentes, valores incompatibles, herramientas no
reconocidas y estados que no pertenecen a las enumeraciones definidas. En este
incremento no se incorpora una base de datos, debido a que el objetivo consiste
en verificar primero el funcionamiento del núcleo y la adecuación de la
estructura común.

\begin{table}[H]
\centering
\caption{Elementos del entorno inicial de implementación}
\label{tab:entorno-implementacion}
\begin{tabularx}{\textwidth}{p{3.5cm}p{3.2cm}X}
\toprule
\textbf{Elemento} & \textbf{Tecnología} & \textbf{Uso dentro del prototipo}\\
\midrule
Lenguaje & Python 3.11 o posterior & Implementación de los modelos, validaciones y procesamiento posterior.\\
Validación de datos & Pydantic 2.x & Validación de tipos, campos obligatorios y valores enumerados.\\
Intercambio de datos & JSON & Representación de entradas, catálogos y datos originales.\\
Pruebas & \texttt{unittest} & Comprobación automática de la estructura y los catálogos.\\
Control de versiones & Git & Registro incremental de los cambios realizados en el artefacto.\\
\bottomrule
\end{tabularx}
\end{table}

\subsection{Arquitectura técnica del prototipo}

La arquitectura técnica se organiza en capas con responsabilidades
diferenciadas. La capa de entrada recibe archivos JSON controlados; la capa de
dominio representa la evaluación, el contexto de ejecución y los resultados
normalizados; la capa de catálogos conserva los elementos configurables del
modelo; y la capa de validación comprueba la consistencia estructural y la
correspondencia entre indicadores y dimensiones.

Esta organización mantiene separada la definición del modelo de la lógica que
posteriormente permitirá importar los reportes producidos por las herramientas.
En consecuencia, los adaptadores de WAVE, axe DevTools y Lighthouse podrán
incorporarse en el siguiente incremento sin modificar la estructura común ya
implementada.

\begin{table}[H]
\centering
\caption{Capas de la arquitectura técnica inicial}
\label{tab:capas-arquitectura-implementacion}
\begin{tabularx}{\textwidth}{p{3.2cm}p{4.3cm}X}
\toprule
\textbf{Capa} & \textbf{Elementos} & \textbf{Responsabilidad}\\
\midrule
Entrada controlada & Archivos JSON de ejemplo & Proporcionar datos representativos para la verificación inicial.\\
Dominio & Evaluación, contexto y resultado normalizado & Representar los elementos principales de información del modelo.\\
Catálogos & Dimensiones, indicadores, estados, prioridades y reglas & Mantener configurables los elementos definidos durante el diseño.\\
Validación & Modelos tipados y comprobaciones cruzadas & Rechazar estructuras inválidas e inconsistencias entre indicador y dimensión.\\
Interfaz inicial & Comando de consola & Ejecutar la validación de uno o más archivos de entrada.\\
\bottomrule
\end{tabularx}
\end{table}

\subsection{Implementación de la estructura común de resultados}

La clase principal del dominio corresponde al resultado normalizado. Sus
atributos mantienen una correspondencia directa con los campos CR-01 a CR-18
definidos durante el diseño. Los campos relacionados con identificadores,
herramienta, tipo de resultado, indicador, dimensión, estado y fecha son
obligatorios. La unidad, ubicación y prioridad conservan la obligatoriedad
condicional establecida en la especificación.

Además de los dieciocho campos, la implementación incorpora el atributo
\texttt{datos\allowbreak Originales}. Este atributo no reemplaza ningún elemento del diseño,
sino que permite cumplir la regla RN-01 mediante la conservación del registro
fuente completo. De esta manera, el resultado normalizado puede ser auditado y
comparado posteriormente con la salida original de la herramienta.

La validación inicial también exige evidencia cuando el registro corresponde a
un hallazgo o una alerta, y exige una unidad cuando corresponde a una métrica.
Estas restricciones permiten detectar entradas incompletas antes de aplicar las
reglas de procesamiento.

\subsection{Implementación de los catálogos}

Los catálogos fueron implementados mediante archivos JSON independientes de la
lógica principal. Esta separación permite modificar o versionar indicadores,
estados y reglas sin alterar las clases del dominio. El contenido inicial se
deriva directamente de la especificación de diseño.

\begin{table}[H]
\centering
\caption{Contenido de los catálogos implementados}
\label{tab:catalogos-implementados}
\begin{tabularx}{\textwidth}{p{4.2cm}p{2.2cm}X}
\toprule
\textbf{Catálogo} & \textbf{Cantidad} & \textbf{Contenido}\\
\midrule
Dimensiones & 3 & Accesibilidad, rendimiento y usabilidad automatizable.\\
Indicadores de accesibilidad & 9 & Indicadores A-01 a A-09.\\
Indicadores de rendimiento & 6 & Indicadores P-01 a P-06.\\
Indicadores de usabilidad automatizable & 5 & Indicadores U-01 a U-05.\\
Estados normalizados & 5 & Estados EN-01 a EN-05.\\
Reglas de normalización & 10 & Reglas RN-01 a RN-10.\\
Prioridades & 5 & Crítica, alta, media, baja y revisión.\\
\bottomrule
\end{tabularx}
\end{table}

El validador comprueba que cada resultado utilice un indicador existente y que
la dimensión declarada corresponda con la dimensión principal del indicador.
Esta comprobación evita, por ejemplo, que un indicador de accesibilidad sea
registrado accidentalmente como resultado de rendimiento.

\subsection{Datos controlados para la verificación inicial}

Para verificar la estructura se prepararon tres archivos de ejemplo. Estos
archivos son simulados y no constituyen resultados de la aplicación del modelo
sobre el Portal Académico. Cada ejemplo representa una forma de dato diferente,
como se resume en el Cuadro~\ref{tab:datos-controlados-implementacion}.

\begin{table}[H]
\centering
\caption{Datos controlados del primer incremento}
\label{tab:datos-controlados-implementacion}
\begin{tabularx}{\textwidth}{p{3.0cm}p{3.4cm}p{2.8cm}X}
\toprule
\textbf{Herramienta} & \textbf{Ejemplo} & \textbf{Indicador} & \textbf{Propósito}\\
\midrule
WAVE & Error de contraste & A-04 & Verificar la representación de un hallazgo con evidencia y prioridad.\\
axe DevTools & Regla \texttt{meta-viewport} & U-02 & Verificar una violación booleana asignada a usabilidad automatizable.\\
Lighthouse & Largest Contentful Paint & P-02 & Verificar una métrica numérica con unidad y estado intermedio.\\
\bottomrule
\end{tabularx}
\end{table}

\subsection{Verificación del primer incremento}

La verificación se realizó mediante pruebas automáticas orientadas a comprobar
la estructura y la coherencia de los catálogos. Se incluyeron casos válidos y
casos deliberadamente incorrectos, como la ausencia de un campo obligatorio,
el uso de una herramienta desconocida, un hallazgo sin evidencia y una
dimensión incompatible con el indicador.

\begin{table}[H]
\centering
\caption{Casos de prueba del primer incremento}
\label{tab:pruebas-primer-incremento}
\begin{tabularx}{\textwidth}{p{1.2cm}X p{3.0cm}}
\toprule
\textbf{ID} & \textbf{Caso de prueba} & \textbf{Resultado esperado}\\
\midrule
CP-01 & Validación de un resultado completo de WAVE. & Aceptado.\\
CP-02 & Omisión del identificador de evaluación. & Rechazado.\\
CP-03 & Uso de una herramienta no contemplada. & Rechazado.\\
CP-04 & Registro de un hallazgo sin evidencia. & Rechazado.\\
CP-05 & Registro de una métrica sin unidad. & Rechazado.\\
CP-06 & Comprobación de las cantidades de los catálogos. & Aceptado.\\
CP-07 & Indicador y dimensión compatibles. & Aceptado.\\
CP-08 & Uso de un indicador inexistente. & Rechazado.\\
CP-09 & Uso de una dimensión incompatible con el indicador. & Rechazado.\\
\bottomrule
\end{tabularx}
\end{table}

Las pruebas confirman que la base del prototipo puede representar resultados de
las tres fuentes previstas y detectar inconsistencias antes del procesamiento.
Con ello se obtiene el primer incremento ejecutable del artefacto. La siguiente
etapa implementará los adaptadores que transformarán los reportes reales de
WAVE, axe DevTools y Lighthouse hacia esta estructura común.
